\documentclass[11pt,letterpaper]{article}
\usepackage[margin=1in]{geometry}
\usepackage[T1]{fontenc}
\usepackage[utf8]{inputenc}
\usepackage{lmodern}
\usepackage{longtable,booktabs,array,enumitem,hyperref,xcolor,graphicx,pdflscape}
\usepackage{indentfirst}
\usepackage{fancyhdr}
\usepackage{graphicx}
\usepackage{titling}
\usepackage{enumitem}
\hypersetup{colorlinks=true,linkcolor=blue,urlcolor=blue}
\setlist[itemize]{leftmargin=*}

% Automatically prefixes the section number to all first-level enumerate lists
% label=\thesubsection.\arabic* pulls "1.1" and adds the list item number
% leftmargin=3\parindent indents 3 paragraphs.
\setlist[enumerate,1]{label=\thesubsection.\arabic*,leftmargin=\parindent}

\pagestyle{fancy}
\geometry{a4paper, margin=1in}

\newcommand{\planned}{\textit{(still to be confirmed)}}
\newcommand{\myTitle}{Golden Eagle Flight Plan Snapshot Objectives}

\begin{document}

\title{Golden Eagle Flight Plan\\Snapshot Objectives\\CS 3338 Final Project}
\author{Project Group 26\\Jason Reese \\ Juan Torres Llamas \\ Victor Hernandez}
\date{November 19, 2026}

\newpage %
\maketitle

\newpage %

\begin{table}[h!]
\centering
\caption{\textbf{Revision History}}
\begin{tabular}{|c|c|c|c|}
\hline
\textbf{Name} & \textbf{Date} & \textbf{Reasons for Changes} & \textbf{Version} \\
\hline
Project Group 26 & 11/17/2026 & First Draft & 1.0.0 \\
\hline
Project Group 26 & 11/19/2026 & Snapshot 1 update & 1.1.0 \\
\hline

\hline

\hline
\end{tabular}
\end{table}

\newpage %

\begin{center}Design draft --- November 2026\end{center}
\tableofcontents
\newpage %

\fancyhf{}
\pagestyle{fancy}
\fancyhead[L]{\myTitle}
\fancyhead[R]{Page \thepage}
\fancyfoot[C]{}

\newpage %



\section{Introduction}
\subsection{Purpose}
This document outlines the objectives for each snapshot of the Golden Eagle Flight Plan project.

There are 4 Snaphots:
	\begin{enumerate}[leftmargin=3\parindent]
		\item Snapshot 1 (Start)
		\item Snapshot 2 (Checkpoint 1)
		\item Snapshot 2 (Checkpoint 3)
		\item Snapshot 4 (Original Due Date)
	\end{enumerate}

\subsection{Document Conventions}
\subsection{Intended Audience and Reading Suggestions}
The intended audience for this document is stakeholders and project management.
This document outlines the targets for each snapshot and the work that will be done.

\newpage %
\section{Overview}
Our goals for the start is to set up the base framework of the project and have a base website. The objective for the first 4 months of our project is to have a full framework for the website with some features implemented including giving different types of accounts (students, faculty, admin) different permissions, and allowing different types of accounts to access certain parts of the website while restricting access to other parts. Everyone on our team will be able to work on both the frontend and backend since they both use JavaScript.

The first Snapshot is fine grained and covers the first month of development towards an MVP (Minimum Viable Product), which is defined as the basic site running with users able to access it.
Snapshots 2-4 have a wider timeline and are not as fine grained and cover features that are being added.


\newpage %
\section{Snapshot 1 (Start)}
Snapshot 1 fine grained and covers the first month of development towards an MVP (Minimum Viable Product), which is defined as the basic site running with users able to access it.

\subsection {Month 1}
In the first month we will be configuring development environments and designing the project.


\subsection {Sprint 1: Define Requirements}
The first week we will build the requirements for the Minimum Viable Product (MVP) of the Golden Eagle Flight plan.
This includes the most basic functionality we require:
	\begin{enumerate}[leftmargin=3\parindent]
		\item User Type requirements.
		\item Login/Logout requirements.
		\item User Creation requirements.
		\item Hosting solution choice.
		\item Technology Stack selection (Frontend/Backend/Database)
		\item Development environments configuration and setup.
	\end{enumerate}

At the end of the first week we will have the requirements for our user types, login/logout and will have determined where we are hosting the site, what technology stack we are using, and have basic development environments setup.

\subsection {Sprint 2: MVP Design}
The second week, we will begin the design for the MVP, which is a site that allows users to login, logout and can be hosted somewhere.
This design will include selection of the technology stack and selection of the hosting provider.
	\begin{enumerate}[leftmargin=3\parindent]
		\item User Creation design.
		\item Login/Logout/User design.
		\item User Type Desgin.
		\item User Database Schema design.
		\item Basic site running in development and deployed.
	\end{enumerate}

At the end of the second week, we will have our site running and engineers will be ready to start development.
We will start with design of User Types, and login/logout, with a stretch goal of having a site that has login and logout functionality and is ready to start adding users to.


\subsection {Sprint 3: MVP Basics}
The 3rd week we will work on the basics of making sure we can create users, login and logout.
This is when the first implementations of roles and approvals will be working.

	\begin{enumerate}[leftmargin=3\parindent]
		\item User Creation implementation.
		\item Login/Logout/User implementation.
		\item Initial Permissions work, and users can be created and login and logout.
		\item Admins can create any type of user.
	\end{enumerate}


\subsection {Sprint 4: Setting Roles}
The 4th week we will create the appropriate roles for users. We will begin assigning permissions to these roles so that users have restricted access based on their roles. Permissions related to specific actions that can be performed on the website will be added when those actions are added.
	\begin{enumerate}[leftmargin=3\parindent]
		\item Creation of Admin, Student, Faculty, and Student Organization roles.
		\item Implementation of a webpage that allows admins to edit permissions. 
		\item Adding global permissions for every logged in user.
		\item Adding role-specific permissions for implemented features.
	\end{enumerate}
\newpage %
\section {Snapshot 2 (Checkpoint 1)}

\subsection {Sprint 5: Homepage}
The first sprint we will work on the design of the homepage the user sees upon logging in. The homepage will contain important announcements and information related to the user which will be added later and it will contain buttons to redirect the user to other parts of the website.
	\begin{enumerate}[leftmargin=3\parindent]
		\item Creation of homepage.
		\item Design of the homepage.
		\item Integrating links into the hoomepage.
		\item Adding a settings dropdown menu to the homepage. 
	\end{enumerate}


\subsection {Sprint 6: }
\newpage %
\section {Snapshot 3 (Checkpoint 3)}
\newpage %
\section {Snapshot 4 (Original Due Date)}

\end{document}
